%! Unit 1: Asking Questions with Data: Study Design — Summative Assessment — Unit Test (blank)
% Shortened 2026-09-29 at the user's request: 20 items, 100 points --- 10 vocabulary, 5 multiple
% choice, 4 one-step computations, 1 critical-thinking short answer. Parallel to ../practice_test/
% item for item.
%
% THE KEY IS GENERATED, NEVER HAND-EDITED:
%   python3 templates/lesson/mkkey.py unit01/tests/actual_test/main.tex \
%       unit01/test_keys/actual_test_key/main.tex unit01/test_keys/actual_test_key/answers.json
% Re-run it after every edit to this file.
\documentclass[10pt]{article}
\usepackage{saar-article}
\usepackage{saar-boxes}

% Work blocks on a test need handwriting room, not typeset spacing. This moves
% the blank and the key together, so it cannot open a page mismatch.
\setlength{\workrowsep}{5pt}

% Part-divider strip
\newcommand{\parthead}[1]{%
  \vspace{6pt}\noindent
  \begin{headlinebox}{cerulean}{\color{white}\bfseries #1}\end{headlinebox}%
  \vspace{4pt}}

% The correct multiple-choice option. Prints nothing in the blank; prints the marker once
% -key is loaded (the same switch that makes `work` blocks visible), so the key regenerates
% from this file byte for byte.
\makeatletter
\newcommand{\correct}{\ifsaar@workvisible\ \textcolor{keyred}{\textbf{$\leftarrow$ correct}}\fi}
\makeatother

\begin{document}

\pageheader{Unit 1: Asking Questions with Data: Study Design}{Unit Test}
\namedateperiod

\vspace{0.06in}
\noindent\textbf{Instructions:} Show all work. No notes unless your teacher says so.
\textbf{100 points:} Part A $20$, Part B $20$, Part C $40$, Part D $20$.

% ======================================================================
\parthead{Part A --- Vocabulary \hfill 10 questions, 2 points each}

\small
Write the letter of the matching description next to each term.

\vspace{2pt}
\begin{minipage}[t]{0.36\linewidth}
\begin{enumerate}[leftmargin=1.9em, itemsep=5pt, topsep=1pt]
  \item \blank{0.8cm} Population
  \item \blank{0.8cm} Sample
  \item \blank{0.8cm} Parameter
  \item \blank{0.8cm} Statistic
  \item \blank{0.8cm} Categorical variable
  \item \blank{0.8cm} Quantitative variable
  \item \blank{0.8cm} Simple random sample
  \item \blank{0.8cm} Convenience sample
  \item \blank{0.8cm} Observational study
  \item \blank{0.8cm} Experiment
\end{enumerate}
\end{minipage}\hfill
\begin{minipage}[t]{0.61\linewidth}
\begin{enumerate}[label=\Alph*., leftmargin=1.9em, itemsep=5pt, topsep=1pt]
  \item The part of the population we actually collect data from.
  \item A variable whose values are measured amounts, like minutes or inches.
  \item The researcher watches or asks, without assigning anyone to a group.
  \item A number that describes a whole population.
  \item Choosing whoever is easiest to reach.
  \item The whole group we want to learn about.
  \item A variable whose values are labels, like a color or a team name.
  \item The researcher assigns people to groups and gives them a treatment.
  \item A number that describes a sample.
  \item Everyone has an equal chance of being picked.
\end{enumerate}
\end{minipage}
\normalsize

% ======================================================================
\boxguard[8]
\parthead{Part B --- Multiple Choice \hfill 5 questions, 4 points each}

\small
\begin{enumerate}[leftmargin=2.2em, itemsep=5pt, topsep=2pt, start=11]

\item Which of these variables is \textbf{quantitative}?
\begin{enumerate}[label=(\Alph*), leftmargin=1.8em, itemsep=1pt, topsep=1pt]
  \item Hair color
  \item Minutes spent reading\correct
  \item Favorite band
  \item Phone area code
\end{enumerate}

% keep each stem with its four options
\boxguard[7]
\item Hillcrest High School has $1{,}200$ students. A survey of $80$ of them finds
that $45\%$ play a sport. The $45\%$ is a
\begin{enumerate}[label=(\Alph*), leftmargin=1.8em, itemsep=1pt, topsep=1pt]
  \item parameter, because it describes the sample.
  \item parameter, because it describes all $1{,}200$ students.
  \item statistic, because it describes all $1{,}200$ students.
  \item statistic, because it describes the sample.\correct
\end{enumerate}

% keep each stem with its four options
\boxguard[7]
\item A coach splits the team roster \textbf{by grade} and then randomly picks $3$
players \textbf{from each grade}. This is a
\begin{enumerate}[label=(\Alph*), leftmargin=1.8em, itemsep=1pt, topsep=1pt]
  \item stratified sample.\correct
  \item systematic sample.
  \item cluster sample.
  \item convenience sample.
\end{enumerate}

% keep each stem with its four options
\boxguard[7]
\item To learn how much town residents read, a researcher surveys people
\textbf{inside the public library}. The results will most likely
\begin{enumerate}[label=(\Alph*), leftmargin=1.8em, itemsep=1pt, topsep=1pt]
  \item be about right, because the people were chosen fairly.
  \item underestimate how much residents read.
  \item overestimate how much residents read.\correct
  \item be right if the researcher surveys more people at the library.
\end{enumerate}

% keep each stem with its four options
\boxguard[7]
\item Researchers draw names from a hat to put $60$ volunteers into two groups. One
group drinks water during a test and the other does not. This study is
\begin{enumerate}[label=(\Alph*), leftmargin=1.8em, itemsep=1pt, topsep=1pt]
  \item an observational study, because the volunteers were watched.
  \item an observational study, because there are two groups.
  \item an experiment, because it measures test scores.
  \item an experiment, because the researchers assigned the groups.\correct
\end{enumerate}

\end{enumerate}
\normalsize

% ======================================================================
\boxguard[12]
\parthead{Part C --- Short Computations \hfill 4 questions, 10 points each}

\small
\begin{enumerate}[leftmargin=2.2em, itemsep=10pt, topsep=2pt, start=16]

% ---------------------------------------------------------------- C1 (1.2)
\item \textbf{Pine Valley Aquatic Center} has $1{,}200$ members. The director asks
$60$ members chosen at random, and $21$ of them take a swim class.

\vspace{2pt}
(a) What percent of the sample takes a swim class?

\begin{work}
  \dfrac{21}{60} &= 0.35 = 35\%
\end{work}

(b) About how many of all $1{,}200$ members take a swim class?

\begin{work}
  0.35 \times 1200 &= 420 \text{ members}
\end{work}

% ---------------------------------------------------------------- C2 (1.3)
\boxguard[14]
\item \textbf{Summit Middle School} has $800$ students and wants a
\textbf{stratified} sample of $100$ --- the same fraction from each grade.

\vspace{2pt}
(a) What fraction of the school is being sampled?

\begin{work}
  \dfrac{100}{800} &= \dfrac{1}{8}
\end{work}

(b) Complete the table.

\vspace{2pt}
\renewcommand{\arraystretch}{1.2}
\begin{tabularx}{\linewidth}{@{} X c c c | c @{}}
\rowcolor{frost}
\textbf{Grade} & 6 & 7 & 8 & \textbf{Total} \\
\midrule
Students in grade & $320$ & $280$ & $200$ & $800$ \\
Number sampled & \blank{1.1cm} & \blank{1.1cm} & \blank{1.1cm} & \blank{1.1cm} \\
\end{tabularx}

% ---------------------------------------------------------------- C3 (1.3)
\boxguard[14]
\item \textbf{Willow Creek Apartments} has $800$ households on a numbered list. The
manager wants a \textbf{systematic} sample of $40$.

\vspace{2pt}
(a) Find the interval $k$.

\begin{work}
  k &= 800 \div 40 = 20
\end{work}

(b) The random start is household $6$. List the first four households chosen.

\vspace{2pt}
\blank{1.6cm} \quad \blank{1.6cm} \quad \blank{1.6cm} \quad \blank{1.6cm}

% ---------------------------------------------------------------- C4 (1.6)
\boxguard[14]
\item In an \textbf{experiment}, $50$ students are randomly assigned to play a
timed math game and $50$ to study the usual way. Afterwards, $35$ of the game group
and $25$ of the usual group passed a quiz.

\vspace{2pt}
(a) Find the percent who passed in each group.

\begin{work}
  \text{Game: } \dfrac{35}{50} &= 0.70 = 70\% \\
  \text{Usual: } \dfrac{25}{50} &= 0.50 = 50\%
\end{work}

(b) How many percentage points apart are the two groups? \blank{2.8cm}

\end{enumerate}
\normalsize

% ======================================================================
\boxguard[12]
\parthead{Part D --- Critical Thinking \hfill 1 question, 20 points}

\small
\begin{enumerate}[leftmargin=2.2em, itemsep=10pt, topsep=2pt, start=20]

\item A random sample of Hillcrest students found that students who \textbf{play a
musical instrument} have higher math grades than students who do not. Nobody was
assigned --- students chose for themselves. The principal says, \emph{``Playing an
instrument causes higher math grades.''}

\vspace{3pt}
Is the principal's conclusion justified? Explain your reasoning, and name one
\textbf{other variable} that could explain the result.

\par\writeline
\par\writeline
\par\writeline
\par\writeline
\par\writeline
\par\writeline

\end{enumerate}
\normalsize

\end{document}
